\section{Alternative Materials}
\subsection{PCB Substrate
\label{subsec:PCBsubstrateAlt}}
\textcolor{red}{\textbf{partial list!} still to add: the RO4003 data sheet values (it was only simulated, in the quick screening runs), price, lead time, and the stock thickness of each material for a 0.381 mm board.}
\begin{itemize}
    \item Rogers RO3003
    \item Rogers TMM10
    \item Rogers TMM10i
\end{itemize}
Values from the Rogers data sheets (10 GHz):
\begin{center}
\small
\begin{tabular}{l|ccc}
\toprule
 & RO3003 & TMM10 & TMM10i \\ \midrule
$\epsilon_r$ (process) & $3.00\pm0.04$ & $9.20\pm0.230$ & $9.80\pm0.245$ \\
$\epsilon_r$ (design) & 3.00 & 9.8 & 9.9 \\
$\tan\delta$ & 0.0010 & 0.0022 & 0.0020 \\
CTE $x/y/z$ (ppm/K) & 17/16/25 & 21/21/20 & 19/19/20 \\
thermal conductivity (W/(m\,K)) & 0.50 & 0.76 & 0.76 \\
\bottomrule
\end{tabular}
\end{center}
Standard thicknesses: RO3003 0.127, 0.254, 0.508, 0.762, 1.524 mm. TMM10 and TMM10i 0.381, 0.635, 0.762, 1.270, 1.524, 1.900, 2.500, 3.175, 3.810, 5.080, 6.350 and 12.70 mm.\par
\noindent What the simulations of this project say, on 0.381 mm boards: the launch alone reflects $|\Gamma| = 0.333$ on RO3003 and $0.467$ on TMM10i; the S21 ripple between the two launches has a period of 3.25 GHz on RO3003 and 2.05 GHz on TMM10i, and is smaller on RO3003; the 50 $\Omega$ trace is 0.40 mm on RO3003 and 0.20 mm on TMM10i. TMM10 was not simulated.
